\section{Introduction}
\label{sec:introduction}

\begin{quote}
\emph{``What are the important problems of your field?\,\ldots\ If you do
not work on an important problem, it's unlikely you'll do important
work.''}\\
\hspace*{\fill}--- Richard Hamming, \emph{You and Your Research} (1986)
\end{quote}

Hamming's question is about asking, not answering. Current AI is
optimized for the opposite: benchmarks, training objectives, and
deployment all reward answering questions someone else has already
posed. Yet the scientific enterprise is bottlenecked earlier in the
loop---someone must notice that a question is worth asking. Recognizing
a productive question---an unexplained discrepancy, a conclusion resting
on a single dataset, a methodological assumption quietly doing
load-bearing work---is a distinct capability from answering it, and
arguably the rarer one. We take the position that this capability,
\emph{scientific question discovery}, is a core requirement for general
scientific intelligence, and that it can be studied rigorously today.

Studying question discovery requires discipline that question answering
does not. An answer can be checked against ground truth; a question's
value only reveals itself later, through the community's engagement with
it. This paper contributes a methodology that makes question discovery
measurable despite that asymmetry:

\begin{enumerate}
  \item \textbf{A framework} (Section~\ref{sec:framework}) that
    represents published evidence as provenance-carrying claims, detects
    and types cross-paper tensions in an evidence graph, refines
    surviving signals into falsifiable questions, and ranks them with a
    two-stage protocol that separates scientific priority from execution
    priority.
  \item \textbf{A corpus discipline} (Section~\ref{sec:benchmark}) in
    which every experiment is bound to a manifest that pins queries,
    versions, and---critically---a time boundary. Data need not live in
    the repository, but the queries, versions, time boundaries, and
    processing code that produce it must.
  \item \textbf{A historical validation protocol}
    (Section~\ref{sec:benchmark}) that exploits the time boundary:
    questions are generated from evidence available before a cutoff and
    judged against the literature published after it, which the system
    never saw. This prevents the system from peeking at the future and
    converts question value into a measurable outcome.
  \item \textbf{An end-to-end instantiation}
    (Section~\ref{sec:experiments}) on exoplanet atmospheres, chosen
    because it uniquely combines papers, structured catalogs (NASA
    Exoplanet Archive), and space-telescope archives (JWST/HST via
    MAST), allowing the full loop---evidence, questions, falsification,
    ranking---to run end to end.
\end{enumerate}

In the historical backtest, every question generated from
pre-2021 evidence was substantively engaged by the 2021--2026 literature;
two were answered, including one whose premise the community later
explicitly refuted---the strongest form of question-discovery success,
since the system had isolated a published conclusion that was in fact an
artifact awaiting correction.
